\section{Conclusion}

We present Nanbeige4.2-3B, a compact model that combines broad agentic capabilities with strong reasoning under a 3B non-embedding parameter budget. Through a Looped Transformer architecture, pre-training from scratch on an improved 28T-token corpus, and scalable SFT trajectory construction and multi-stage RL recipes, Nanbeige4.2-3B achieves solid performance across code-agent, office-agent, and complex tool-use tasks while remaining highly competitive in mathematical, coding, and scientific reasoning. Its performance in local personal-assistant workflows further demonstrates the practical potential of compact agentic models.

Looking forward, we will continue improving model capacity and efficiency under a constrained computation budget. On the architecture side, building on Nanbeige4.2, we are improving how layers are organized within the loop, how information is passed and reused across different depths, and how token embeddings incorporate local n-gram context as external memory.\footnote{Preliminary implementations are already included in the modeling code released with Nanbeige4.2.} These directions aim to improve model performance with minimal additional FLOPs. We are also exploring linear and sparse attention mechanisms for more efficient long-context modeling. On the agent side, we will continue scaling post-training along three dimensions: task difficulty, the number and diversity of environments, and trajectory quality. Other than post-training, we plan to study how agentic capabilities can be established earlier during pre-training. We will train and release larger Nanbeige models and investigate how to use them to distill a stronger small model.
